% source-figure:q09-flowchart, fixed science PDF p.1 / printed 961.
% Uses only basic TikZ paths and nodes; scaling/spacing are presentational.
\begin{tikzpicture}[>=stealth, every node/.style={font=\small}]
  \draw[rounded corners=5pt] (-0.52,0.05) rectangle (0.52,0.55);
  \node at (0,0.30) {开始};
  \draw[->] (0,0.05) -- (0,-0.45);
  \draw (-0.62,-0.45) -- (0.82,-0.45) -- (0.62,-0.95) -- (-0.82,-0.95) -- cycle;
  \node at (0,-0.70) {输入 $t$};
  \draw[->] (0,-0.95) -- (0,-1.45);
  \draw (-1.70,-1.45) rectangle (1.70,-2.00);
  \node at (0,-1.725) {$S=1,\ n=0,\ m=\frac12$};
  \draw[->] (0,-2.00) -- (0,-2.50);
  \draw (-0.92,-2.50) rectangle (0.92,-3.00);
  \node at (0,-2.75) {$S=S-m$};
  \draw[->] (0,-3.00) -- (0,-3.48);
  \draw (-1.50,-3.48) rectangle (1.50,-4.03);
  \node at (0,-3.755) {$m=\frac m2,\ n=n+1$};
  \draw[->] (0,-4.03) -- (0,-4.53);
  \draw (-1.08,-4.85) -- (0,-4.53) -- (1.08,-4.85) -- (0,-5.17) -- cycle;
  \node at (0,-4.85) {$S>t$};
  \draw[->] (-1.08,-4.85) -- (-2.05,-4.85) -- (-2.05,-2.25) -- (0,-2.25) -- (0,-2.50);
  \node[left] at (-1.08,-4.72) {是};
  \draw[->] (0,-5.17) -- (0,-5.65);
  \node[right] at (0.06,-5.40) {否};
  \draw (-0.62,-5.65) -- (0.82,-5.65) -- (0.62,-6.15) -- (-0.82,-6.15) -- cycle;
  \node at (0,-5.90) {输出 $n$};
  \draw[->] (0,-6.15) -- (0,-6.62);
  \draw[rounded corners=5pt] (-0.52,-7.12) rectangle (0.52,-6.62);
  \node at (0,-6.87) {结束};
\end{tikzpicture}
